\section{Introduction}
\label{sec:intro}

Large language models (LLMs) are increasingly deployed under strict memory and
latency budgets. Post-training quantization (PTQ) maps FP16 weights to 2--4
bits without retraining, but aggressive quantization---especially at 2
bits---can collapse performance on structured tasks such as text-to-SQL, math
reasoning, and code generation. Mixed-precision PTQ mitigates this by keeping a
small set of outlier weights in FP16 while quantizing the
rest~\citep{dettmers2022llmint8,frantar2023gptq,dettmers2023spqr}.

\paragraph{The core problem.}
Outlier selection is a \emph{localization} problem: which weights matter for a
downstream task once quantization noise is injected? Uniform salience
heuristics (magnitude, activation scale) ignore task structure. Calibration on
generic corpora misses task-specific circuits. Even strong mixed-precision
methods struggle in the 2-bit regime, where the allowed FP16 budget (often
$<$1\% of weights) must be spent with surgical precision.

\paragraph{Polysemantic weights vs.\ monosemantic features.}
A separate challenge is \emph{where} in the network task signal lives.
Individual MLP neurons and their connecting weights are typically
polysemantic---each participates in many unrelated computations via
superposition~\citep{elhage2021mathematical,bricken2023monosemanticity}. A
gradient on cross-entropy therefore reflects a mixture of all roles a weight
plays across the vocabulary and prompt positions, not a single task-specific
mechanism. \baseline{} inherits this ambiguity: although it conditions on task
calibration data, saliency is still computed directly on physical weights from
CE gradients, so the selected outliers may protect broadly useful directions
rather than the sparse features that implement SQL grounding, aggregation, or
schema linking.

\paragraph{Task-Circuit Quantization (\baseline{}).}
\citet{xiao2025tacq} frame outlier selection as automated circuit discovery for
PTQ. They score each weight by combining (i) contrastive gradients on task
calibration data, (ii) expected weight change under quantization
$\Delta W = W - W_{\text{quant}}$, and (iii) a global top-$p$ mask under a
fixed memory budget. Crucially, \baseline{} operates in weight space on
polysemantic gradient signal; it does not decompose MLP computation into
features before ranking outliers.

\paragraph{Step 1 --- \method{}: route saliency through a monosemantic basis.}
Our starting point is interpretability-guided: if CE gradients on weights are
polysemantic, attribute through a basis that is not. \method{}
(Transcoder-guided Dynamic Sparse Overlay) uses sparse TopK
transcoders~\citep{dunefsky2024transcoders} to identify features that activate
preferentially on task-faithful versus corrupted inputs, back-propagates a
feature-filtered reconstruction objective, and \emph{intersects} the resulting
circuit saliency with standard CE saliency (geometric mean;
\S\ref{sec:method}). This works: +2.9 points mean Spider execution accuracy
over four paired seeds and +3.3 points GSM8K at 2-bit versus \baseline{} at
matched budget, with the selected mask sharing only $J{=}0.30$ of its weights
with \baseline{}'s.

\paragraph{Step 2 --- dissect it: which ingredient carries the gain?}
A method with three coupled ingredients (contrastive circuit localization, a
pretrained transcoder dictionary, and a conjunction rule) invites the question
of what is actually doing the work---and transcoders exist only for a handful
of models, which would cap the method's reach. We therefore run a
\emph{saliency-source ablation} in which everything is held fixed (model,
0.35\% budget, contrastive pairs, GPTQ, evaluation) and only the signal that
ranks weights varies across eight arms: random, weight magnitude $|W|$,
magnitude $\times$ quantization error $|W|\cdot|\Delta W|$, CE
($\approx$\baseline{}), circuit-only, and the conjunction---each circuit arm
instantiated in \emph{two} bases: the transcoder and, replacing it, the
model's \emph{own MLP neurons} with \texttt{down\_proj} as the decoder.

\paragraph{Step 3 --- the distilled method: \methodfree{}.}
The ablation delivers a clean dissection: the conjunction is the active
ingredient, and the dictionary is not. The native-basis conjunction
(\methodfree{}) recovers most of \method{}'s gain on GSM8K and surpasses it on
MMLU-STEM, while requiring no external dictionary, SAE, or transcoder---the
method applies out-of-the-box to any transformer. What began as a
transcoder-dependent technique ends as a general principle: \emph{protect a
weight only if it is salient to both the task loss and the task circuit; any
sufficiently task-discriminative feature basis---including the model's
own---suffices to define the circuit.}

\paragraph{Contributions.}
\begin{itemize}
  \item \textbf{\method{}:} we introduce transcoder-guided outlier selection
        for PTQ---routing saliency through task-discriminative monosemantic
        features and intersecting with CE---and show it beats \baseline{} at
        matched budget on Spider (+2.9\,pp mean over four paired seeds) and
        GSM8K 2-bit (+3.3\,pp), with interpretability analysis showing low
        mask overlap ($J{=}0.30$) and causal knockout ($-29$\,pt).
  \item \textbf{Dissection:} a controlled eight-arm saliency-source ablation
        isolates the active ingredient. Task-circuit information dominates
        magnitude/error/random placement by 10--26 points at 2-bit, and the
        \emph{conjunction} beats both of its parts (GSM8K: 28.7\% vs.\ CE
        25.4\% / circuit-only 25.6\%; \S\ref{sec:ablation}).
  \item \textbf{\methodfree{}:} the dissection shows the dictionary is
        optional. Using the model's own MLP neurons as the circuit basis
        recovers most of the transcoder gain on GSM8K and \emph{surpasses} it
        on MMLU-STEM (44.2\% vs.\ 38.9\%), removing the dependency on external
        interpretability artifacts and extending task-circuit PTQ to any
        transformer.
  \item \textbf{Regime analysis:} all task-conditioned methods converge at
        3-bit; conjunctive saliency is a \emph{low-bit} phenomenon
        (\S\ref{sec:results-3bit}), locating exactly where saliency choice
        matters.
  \item We release code, contrastive calibration sets, and reproduction
        scripts for \baseline{}, \method{}, and \methodfree{} masks.
        \todo{anonymize for submission}
\end{itemize}

\paragraph{Paper roadmap.}
Section~\ref{sec:encoded-used} gives the theoretical framing;
Section~\ref{sec:rqs} states research questions;
Section~\ref{sec:related} situates our work;
Section~\ref{sec:method} describes the algorithm and the dictionary-free
variant; Sections~\ref{sec:setup}--\ref{sec:results} cover setup and findings.
